\section{The Standard Model}

The matter and force of the world are not added to the model. They are read off the substrate's own furniture, the dock with its $24$ sites and the ternary tone each site carries. This section shows how that small base builds a grand-unified gauge structure, how the structure breaks down to the gauge group of the Standard Model, and how one full family of quarks and leptons falls out as a single object. It also marks, plainly, the two places where the substrate stops short and a further principle would be needed.

Two ideas recur, so a plain statement of each helps before the formal work begins. A \emph{gauge group} is the symmetry that governs how a force acts. The Standard Model has the gauge group $\mathrm{SU}(3) \times \mathrm{SU}(2) \times \mathrm{U}(1)$, one factor for the strong force, one for the weak, one for electromagnetism. A \emph{generation} is one complete set of the matter particles, the up and down quarks, the electron, and its neutrino, in the lightest case. There are three generations in nature, each a heavier copy of the same pattern. The claim of this section is that the dock and the tone build the gauge group and one generation together, as a single algebraic fact.

\begin{principle}[The gauge group and one family come from the dock plus the tone]
The $24$ sites of a dock carry the $D_4$ root system, which is the symmetry $\mathrm{so}(8)$. Adding the ternary tone as one more axis promotes it to $\mathrm{so}(10)$. The $16$-component spinor of $\mathrm{so}(10)$ is exactly one generation of matter. A self acting as a condensate breaks $\mathrm{so}(10)$ down to the Standard Model. Nothing in this chain is added by hand.
\end{principle}

\subsection{From the sites to \texorpdfstring{$\mathrm{so}(8)$}{so(8)}}

Recall the substrate. The mesh is the order-$4$ honeycomb $\{3,4,3,4\}$ of flat four-dimensional space, a weave of docks. Each dock is one cell, a here-now, and out of each dock run exactly $24$ sites, the directed slots where a tone can live. There is no twenty-fifth site and no rest slot. These $24$ sites are not a loose collection. They are the $24$ vertices of the $24$-cell, equivalently the $24$ roots of the $D_4$ root system, and a root system is the fingerprint of a Lie algebra. The Lie algebra of $D_4$ is $\mathrm{so}(8)$, the rotations of eight-dimensional space.

\begin{definition}[The site algebra]
The $24$ sites of a dock are the roots of $D_4$. The Lie algebra they generate is $\mathrm{so}(8)$, a rank-four algebra with $28$ generators. This is the \emph{site algebra}, the symmetry that the bare dock carries before any further structure is added.
\end{definition}

The site algebra has a special feature, its triality. Under the outer symmetry $S_3$ of order six, the three eight-dimensional representations of $\mathrm{so}(8)$, the vector $8v$ and the two spinors $8s$ and $8c$, are permuted into one another. The $24$ sites decompose as $8v + 8s + 8c$, the vector together with its two spinors. The $8v$ part is the photon sector. The $8s$ and $8c$ parts are the two chiralities of the fermion. This decomposition is the foundation of the matter sector, and it is the reason the substrate carries spin at all, a point developed where spin and fermions are taken up directly.

\subsection{Why the tone is forced}

The site algebra alone cannot hold the Standard Model. This is not a failure of imagination, it is a theorem about root systems, and it is worth stating with care because it is exactly what forces the tone into the picture.

The Standard Model needs $\mathrm{SU}(3) \times \mathrm{SU}(2)$ to sit inside the gauge algebra, with the two factors acting independently. For two subgroups to act independently inside some $\mathrm{so}(N)$, their root systems must be orthogonal. The $\mathrm{SU}(2)$ factor is generated by a single $A_1$ root, and the $\mathrm{SU}(3)$ factor by an $A_2$ pair, so what is needed is a root orthogonal to a full $A_2$ system.

\begin{lemma}[The orthogonality obstruction]
Every root of $D_4$ has exactly two non-zero entries, of the form $(\pm 1, \pm 1, 0, 0)$ up to permutation. No such vector can be orthogonal to a full $A_2$ root system, because an $A_2$ system already spans a two-dimensional subspace and every $D_4$ root meets it. Therefore $\mathrm{SU}(3) \times \mathrm{SU}(2)$ does not embed orthogonally in $\mathrm{so}(8)$, and the Standard Model does not fit inside the site algebra.
\end{lemma}

The fix is exact. Add the ternary tone as a fifth axis. The tone is the vibe value each site carries, ranging over $\{-1, 0, +1\}$ as pain, peace, and pleasure, and treating its sign as an extra weight extends the rank-four $D_4$ to the rank-five $D_5$. The root system $D_5$ is $\mathrm{so}(10)$, and $\mathrm{so}(10)$ has room for the orthogonal $A_2$ the obstruction demanded. The book-keeping is clean.

\begin{table}[ht]
\centering
\begin{tabular}{@{}lll@{}}
\toprule
structure & roots & meaning \\
\midrule
$D_4$ (the sites) & $24$ & the bare dock, $\mathrm{so}(8)$ \\
the extra & $16$ & $8v$ tensored with the tone sign \\
$D_5 = \mathrm{so}(10)$ & $40$ & the dock plus the tone \\
\bottomrule
\end{tabular}
\caption{Root counting for the promotion from $\mathrm{so}(8)$ to $\mathrm{so}(10)$. The $16$ new roots are the vector $8v$ paired with each of the two tone signs.}
\label{tab:so10-roots}
\end{table}

The point to hold onto is that the tone is not a convenient extra knob. It is forced. The site algebra provably cannot carry the Standard Model, and the smallest fix that the obstruction permits is exactly the one axis the theory already had, the ternary tone.

\begin{result}[The tone promotes the symmetry to \texorpdfstring{$\mathrm{so}(10)$}{so(10)}]
The $24$ sites give $\mathrm{so}(8)$, which cannot hold the Standard Model by the orthogonality obstruction. Adding the ternary tone as a fifth weight builds the $D_5$ root system, which is $\mathrm{so}(10)$. The promotion is forced by the obstruction, not chosen.
\end{result}

\subsection{One generation in one spinor}

Grand unification on $\mathrm{so}(10)$ has a famous virtue. Its smallest spinor representation is $16$-dimensional, and that $16$ is precisely the particle content of one generation, with not one slot too many and not one too few. The substrate produces this $16$ directly, from the two spinors of the site algebra and the tone sign.

\begin{table}[ht]
\centering
\begin{tabular}{@{}lll@{}}
\toprule
half-spinor & tone & content \\
\midrule
$8s$ & $+$ & half of one generation \\
$8c$ & $-$ & the other half \\
\bottomrule
\end{tabular}
\caption{The two half-spinors of the dock, tagged by the sign of the tone, assemble into the $16$ of $\mathrm{so}(10)$.}
\label{tab:so10-spinor}
\end{table}

So $16 = 8s + 8c$, and that $16$ is one full family of matter. The two chiralities that triality gave us, paired with the two signs of the tone, are exactly the spinor of the unified group. The matter content is not assembled by hand from a list of particles. It is one irreducible object that the dock and the tone build together.

\begin{proposition}[The $16$ is one generation]
The $16$-component spinor of $\mathrm{so}(10)$ decomposes as $8s + 8c$ under the site algebra, where $8s$ and $8c$ are the two half-spinors of the dock tagged by the sign of the tone. This $16$ is exactly one generation of Standard Model matter, quarks and leptons together with a right-handed neutrino.
\end{proposition}

\subsection{Breaking to the Standard Model}

A grand-unified symmetry is too large to be the world we see. Something must break it down to $\mathrm{SU}(3) \times \mathrm{SU}(2) \times \mathrm{U}(1)$. In the substrate that something is a self acting as a condensate. A self is a bound, self-correcting, identity-bearing knot of tone woven across many docks over many beats. When such a knot condenses it plays the role of the grand-unified Higgs field, and its presence selects which part of the symmetry survives.

\begin{result}[Breaking by a self-condensate]
A self-condensate breaks $\mathrm{so}(10)$ along the Georgi-Glashow chain,
\[
\mathrm{so}(10) \;\longrightarrow\; \mathrm{su}(5) \;\longrightarrow\; \mathrm{SU}(3) \times \mathrm{SU}(2) \times \mathrm{U}(1).
\]
Any self-condensate leaves $\mathrm{su}(5)$ unbroken, with $20$ surviving roots. The breaking is automatic, not fine-tuned.
\end{result}

Under the intermediate $\mathrm{su}(5)$ the spinor splits into three pieces, and the split is the textbook Georgi-Glashow assignment.

\begin{table}[ht]
\centering
\begin{tabular}{@{}ll@{}}
\toprule
$\mathrm{su}(5)$ piece & content \\
\midrule
$1$ & right-handed neutrino \\
$\overline{5}$ & half the generation \\
$10$ & the rest \\
\bottomrule
\end{tabular}
\caption{The $16$ of $\mathrm{so}(10)$ under $\mathrm{su}(5)$, the complete matter content of one generation.}
\label{tab:su5-split}
\end{table}

So $16 = 1 + 10 + \overline{5}$, the complete matter content of one generation, quarks, leptons, and a right-handed neutrino, with nothing left over.

There is a clean conceptual payoff here. The field that breaks the symmetry is itself a self, a piece of matter. The matter sector and the symmetry-breaking sector are the same object, not two.

\begin{principle}[The self is the Higgs]
The breaking field is a self-condensate. The thing that breaks the grand-unified symmetry is therefore a knot of matter, not a separate scalar grafted on for the purpose. Matter and the Higgs are one sector.
\end{principle}

\subsection{The Weinberg angle}

The first sharp number the construction delivers is the Weinberg angle, the parameter $\theta_W$ that mixes the weak and electromagnetic forces. At the unification scale it is fixed by a sum over the matter of one generation, the ratio of the total squared weak-isospin to the total squared electric charge.

\[
\sin^2 \theta_W = \frac{\sum T_3^2}{\sum Q^2} = \frac{3}{8} \quad \text{exactly at unification.}
\]

This value of $3/8$ is the classic grand-unified prediction, and it follows here from the same group theory that fixed the $16$. It is not the value measured in the laboratory, because the measured value lives at low energy and the prediction lives at the unification scale. The two are joined by the renormalization-group flow, the standard running of couplings with energy.

\begin{table}[ht]
\centering
\begin{tabular}{@{}ll@{}}
\toprule
spectrum & $\sin^2 \theta_W$ at $M_Z$ \\
\midrule
measured & $0.231$ \\
MSSM-like running & reaches $0.231$ \\
bare Standard Model running & about $0.21$ \\
\bottomrule
\end{tabular}
\caption{The Weinberg angle run down from $3/8$ at unification to the $Z$ mass. The supersymmetric-like spectrum reaches the measured value, the bare Standard Model spectrum undershoots.}
\label{tab:weinberg}
\end{table}

Run down with the supersymmetric-like spectrum, the prediction reaches the measured $0.231$, and the three gauge couplings unify cleanly at about $10^{16}$ GeV. So the prediction is confirmed against experiment, and as a bonus it favours low-energy supersymmetry.

\begin{result}[The Weinberg angle]
The substrate predicts $\sin^2 \theta_W = 3/8$ at the unification scale, from the charge content of the $16$. Run down to the $Z$ mass with a supersymmetric-like spectrum it reaches the measured value $0.231$, with the three couplings unifying at about $10^{16}$ GeV.
\end{result}

\subsection{The mass relations}

A few further relations fall out at the unification scale, where the grand-unified symmetry still relates particles that are distinct at low energy. They are stated here with their honest status, since they range from clean to model-dependent.

\begin{table}[ht]
\centering
\begin{tabular}{@{}lll@{}}
\toprule
relation & statement & status \\
\midrule
$b$-$\tau$ unification & $m_b = m_\tau$ at unification, from the shared Yukawa & runs to $m_b/m_\tau \approx 2.3$ at $M_Z$, observed $2.35$ \\
determinant relation & $\det(M_{\ell}) = \det(M_d)$ & from the discrete fact $\operatorname{Tr} Y = 0$ over the $16$ \\
Georgi-Jarlskog & factor-three relations & weakest, model-dependent \\
\bottomrule
\end{tabular}
\caption{Mass relations at the unification scale. The $b$-$\tau$ relation runs to within a few percent of the observed ratio.}
\label{tab:mass-relations}
\end{table}

The $b$-$\tau$ relation is the strongest. A shared Yukawa coupling sets the bottom quark and the tau lepton equal at unification, and running that equality down to the $Z$ mass gives a ratio of about $2.3$, against the observed $2.35$. The determinant relation follows from the discrete fact that the trace of the Yukawa matrix vanishes over the $16$. The Georgi-Jarlskog factor-three relations are the weakest, since they depend on extra model choices the bare substrate does not fix.

\subsection{The gauge field is genuine}

It is one thing to find the right gauge algebra in the symmetry of the dock, and another to show that the substrate actually carries a gauge field. The non-abelian field was verified directly with link matrices, the discrete connection that lives on the sites between docks, in the work on gauge and the photon. The underlying fact is structural.

\begin{principle}[Gauging is the symmetry of the knit]
A non-abelian gauge group is nothing more than a local symmetry of the knit, the collide-then-stream law. Gauging $\mathrm{so}(10)$ amounts to the existence of an $\mathrm{so}(10)$-symmetric collision, which the substrate supplies. That symmetric collision is forced in the infrared by triality.
\end{principle}

The $\mathrm{U}(1)$ piece of the gauge group is the local conservation of the charge carried by the knit, the same charge that the lean orders. The non-abelian pieces are the symmetric collision. Both are read off the law rather than postulated, and both are taken up in detail where gauge and the photon are treated.

\subsection{What is derived and what is open}

The construction is strong but not complete, and it is important to mark the line exactly. The gauge group, one generation, the breaking, and the Weinberg angle are derived. Two things are not, the value of the gauge coupling and the existence of three distinct generations.

\begin{table}[ht]
\centering
\begin{tabular}{@{}ll@{}}
\toprule
item & status \\
\midrule
$\mathrm{so}(10)$ from the sites plus the tone & derived \\
the tone forced by the orthogonality obstruction & derived \\
$16 =$ one generation & derived \\
breaking to the Standard Model by a self-condensate & derived, automatic \\
$\sin^2 \theta_W = 3/8$ at unification & derived, runs to the measured value \\
the mass relations & derived at unification \\
the coupling value & open \\
three distinct generations & partial \\
\bottomrule
\end{tabular}
\caption{The status of each result in the matter and force sector.}
\label{tab:sm-status}
\end{table}

The coupling value is treated honestly elsewhere. Scanning the gauge coupling, the field a moving charge radiates follows a clean square law in the coupling with no special or critical value anywhere, vanishing only at zero coupling. A pure power law with no distinguished point means the coupling enters as an overall scale, so nothing in the bare knit fixes its magnitude.

\begin{principle}[The knit fixes the form, not the value]
The knit determines the form of the dependence on the gauge coupling, a square law, but not its value. Fixing the value would need a further principle the bare knit does not contain, a renormalization-group fixed point, anomaly matching, or the full charged-particle spectrum with its running. The substrate does not derive $1/137$ and does not claim to.
\end{principle}

The generation count is the program's one genuinely novel bet, that triality seeds three families. The naive reading fails, the triple $8v$, $8s$, $8c$ is a vector plus two chiralities of one generation, not three. The deeper path is the substrate's exceptional symmetry, and it has been computed rather than asserted. The $24$ sites are the long roots of $F_4$, and $F_4$ is the automorphism group of the exceptional Jordan algebra $J_3(\mathbb{O})$, the $27$-dimensional Albert algebra. Its rank-three frame is three primitive orthogonal idempotents summing to the identity, and the three is forced, because the Jordan identity holds for Hermitian octonion matrices at size three and below but fails at size four, the octonions being non-associative.

\begin{result}[A real family symmetry, but degenerate slots]
The substrate forces a rank-three exceptional Jordan structure, with an exact $S_3$ family symmetry permuting the three slots by genuine Jordan automorphisms. But because the $S_3$ is exact, the three slots are degenerate, every quantity the algebra supplies is equal across them. The substrate gives a family symmetry, not a family breaking. Splitting the degeneracy into three distinct masses is Boyle's open conjecture, and it stays open.
\end{result}

\subsection{Why it matters}

The whole gauge and matter sector, the prediction $\sin^2 \theta_W = 3/8$, and the mass relations come from the dock plus the tone, two pieces the theory already had for entirely independent reasons. The Standard Model emerges by an automatic breaking through a self-condensate that is itself a piece of matter, so the Higgs is not an extra ingredient but the matter sector seen from another side. The two endpoints the substrate cannot reach, the value of the coupling and the splitting of the three slots, are pushed to their exact frontier and stated plainly rather than papered over. The gauge group and one full family of quarks and leptons are not assumptions of the model. They are the symmetry of the dock and the tone, read off in order.
